\chapter{多模态大模型入门}\label{多模态大模型入门}

小助已经能够阅读文字和调用工具。现在学生上传一张几何题照片，希望它读出条件、理解图形并给出解答。完成这个任务需要把图像中的文字、位置和形状转成模型能够利用的表示，再结合问题生成回答。多模态大模型把图像、声音、视频等信息纳入模型输入或输出，使助手能够围绕多种信息形式完成任务。本章以图像输入、文字输出的视觉语言模型为例，介绍基本结构与后训练流程。

\section{基本结构 编码器、连接器与语言模型}

一种常见方案由三个部分组成：视觉编码器提取图像特征，模态连接器把视觉特征映射到语言模型使用的表示空间，语言模型结合视觉表示和文字问题生成答案。LLaVA 是这一组合方式的代表工作，后续模型在分辨率、特征压缩和训练数据等方面继续扩展。

\begin{enumerate}
    \item \textbf{视觉编码器。}输入图片被划分为图像块（patch），编码器根据各图像块及其空间关系产生一组视觉向量。预训练视觉编码器提供对形状、物体和图文关联的基础表示。
    \item \textbf{模态连接器。}视觉向量与语言模型的文本嵌入在维度和表示分布上存在差异。线性投影、多层感知机或查询式压缩模块把视觉特征转换为语言模型可以接收的向量序列。
    \item \textbf{语言模型。}将转换后的视觉向量与问题的文本嵌入按模板组织，语言模型通过注意力读取这些输入，逐 token 生成回答。图片占位符用于标记视觉输入插入的位置，实际的图片内容由对应视觉向量提供。
\end{enumerate}

把图片记为 $I$、文字问题记为 $x$，这条路径可以简写为
\[
V=E_{\psi}(I),\qquad U=P_{\omega}(V),\qquad
p_{\theta,\psi,\omega}(y\mid I,x)
=\prod_{t=1}^{T}p_{\theta}(y_t\mid y_{<t},x,U).
\]
这里，$E_{\psi}$ 是参数为 $\psi$ 的视觉编码器，$V$ 是编码器输出的视觉特征序列；$P_{\omega}$ 是参数为 $\omega$ 的模态连接器，$U$ 是转换后送入语言模型的视觉向量序列；$\theta$ 是语言模型参数。$y$ 是回答，$T$ 是回答 token 数，$y_t$ 是第 $t$ 个回答 token，$y_{<t}$ 是此前的回答前缀，乘号表示把各步条件概率相乘。这个公式对应将视觉向量作为输入序列的方案；另一类结构通过跨模态注意力，让文本状态读取单独保存的视觉特征。

\subsection{经典方法 CLIP 怎样学会图文匹配}

视觉编码器的基础能力可以通过图文配对数据学习。2021 年提出的对比语言图像预训练（Contrastive Language--Image Pre-training，CLIP）采用图像编码器与文本编码器组成的双编码器结构，把图片与描述映射到共同的向量空间。训练推动匹配图文的向量相似度升高，使模型能够通过自然语言描述识别视觉内容。

例如，一个批次包含三对数据：猫的图片与``一只猫''，狗的图片与``一只狗''，汽车图片与``一辆汽车''。图像编码器输出三个图像向量，文本编码器输出三个文本向量，两两计算得到一个 $3\times3$ 的相似度矩阵。对每张图片，训练要求对应描述的得分高于另外两条描述；对每条描述，也要求对应图片的得分高于另外两张图片。配对关系由训练数据提供，两个编码器通过对比损失一起更新。

更一般地，对包含 $B$ 对图文的批次，把归一化后的图像向量记为 $u_i$，文本向量记为 $v_j$，可以写出图像到文本的损失：
\[
s_{ij}=\frac{u_i^{\mathsf T}v_j}{\tau},\qquad
L_{I\to T}=-\frac1B\sum_{i=1}^{B}
\log\frac{e^{s_{ii}}}{\sum_{j=1}^{B}e^{s_{ij}}}.
\]
这里，$B$ 是批次中的图文对数，$i$ 是图片编号，$j$ 是文本编号，同编号的图片与文本是一对。$u_i$、$v_j$ 是长度归一化为 1 的向量，转置符号 $\mathsf T$ 用于计算点积，因此 $u_i^{\mathsf T}v_j$ 等于余弦相似度；$\tau>0$ 是温度参数，控制得分差异进入 softmax 后的尺度。$s_{ij}$ 是缩放后的相似度，$\log$ 是自然对数。分式表示在批次的全部描述中选中匹配描述的概率。文本到图像的损失 $L_{T\to I}$ 则固定文本、对图片候选归一化，CLIP 对两个方向的损失取平均。

训练完成后，可以把类别写成``一张猫的照片''、``一张狗的照片''等描述，与待识别图片分别计算相似度，选择得分最高的类别。这种方式实现了用文字指定类别的零样本分类，也可用于图文检索。CLIP 学习的是图文匹配表示；多模态助手进一步接入生成式语言模型，学习围绕图片进行问答。

在前述多模态结构中，可以使用 CLIP 预训练的图像编码器作为视觉起点，取图像块的中间特征，经连接器送入语言模型。用于 CLIP 图文匹配的是整图表示，助手读取的则可以是一组保留空间细节的视觉特征。后续的图文对齐与指令微调，让语言模型学会利用这些特征生成答案；CLIP 的文本编码器完成预训练时的配对任务，生成回答由接入的语言模型负责。

\section{训练路线 从图文对齐到多模态后训练}

\textbf{先学习视觉表示与语言的对应关系。}图文对齐阶段使用图片及其描述，让连接器学习如何把视觉特征交给语言模型。一个易于理解的起点是冻结已有视觉编码器和语言模型，只训练连接器：输入图片，预测对应的文字描述，用文本预测误差调整连接器参数。研究者也可以根据资源与数据规模，解冻部分或全部视觉和语言参数，进行更充分的多模态联合训练。

\textbf{再学习根据图像执行指令。}多模态监督微调使用``图片、问题、目标回答''组成的样本，例如识别物体、读取表格、解释图表和分析几何条件。与第四章的 SFT 相同，训练以给定上下文预测目标回答；上下文现在同时包含视觉表示和文本。一个样本的回答损失可以写成
\[
L_{\mathrm{MM\text{-}SFT}}
=-\sum_{t=1}^{T}\log p_{\theta,\psi,\omega}
(y_t^*\mid I,x,y_{<t}^*).
\]
$L_{\mathrm{MM\text{-}SFT}}$ 表示多模态监督微调损失，$I$ 是图片，$x$ 是用户问题，$y_t^*$ 是目标回答在第 $t$ 步的 token，$y_{<t}^*$ 是目标回答前缀，$T$ 是目标回答长度，$\log$ 是自然对数。公式对回答位置计算损失，图片与问题作为条件输入；具体更新哪些参数，由冻结与解冻配置决定。

\textbf{最后优化回答质量与任务结果。}对于同一图片和问题，可以标注偏好回答与偏恶回答，采用 DPO 等方法训练；第五章 DPO 公式中的条件 $x$ 相应扩展为图片与问题的组合。对于答案可以检查的图表数值题、数学题或界面操作任务，也可以通过结果奖励进行强化学习。数据需要覆盖视觉证据的使用方式：读对对象、对应正确位置、保留单位，再完成推理和回答。

对于以文本后训练为起点的读者，一条直接的实践路线是选择已有多模态模型，准备领域图片与指令数据，先进行小规模 SFT，再按任务需要增加偏好或结果反馈。这样可以把实验重点放在目标行为、数据质量和评估上。

\section{一个例子 区分读图、推理与回答}

学生上传一张柱状图，询问``乙地区的销售额比甲地区高百分之几？''。图中甲地区为 100 万元、乙地区为 125 万元。一次完整求解包含三个环节：视觉部分识别柱子、标签和数值；语言模型把问题对应到 $(125-100)/100$；最终回答组织为``乙地区比甲地区高 25\%''，并说明所用数据。

准备后训练样本时，可以围绕同一图片设置不同问题：先问各地区销售额，再问差额，最后问增长比例。这样既覆盖视觉读取，也覆盖计算与表达。数据中保留数值、单位、标签和计算口径，便于检查模型在哪个环节发生变化。

评估时也可以分层测量。先让模型单独输出图中数值，检查读图准确率；再把正确数值写成文字问题，检查计算能力；最后使用原图完成端到端问答，比较完整任务表现。若文字条件下计算正确、图片条件下数值读取错误，可以优先调整分辨率、视觉训练数据或图表读取样本。

\section{图片、视频与语音 怎样扩展输入}

\textbf{高分辨率与多图输入。}文档小字和细密图表需要保留局部细节。动态分辨率、图像分块与局部裁剪把更多视觉信息送入模型，同时增加视觉向量数量。后训练的数据处理应保持图片、裁剪位置与问题对应，并按视觉输入长度安排批次与显存预算。多图任务则需要保留图片顺序和标识，让模型能够比较不同图片中的对象。

\textbf{视频输入。}视频可以看作带有时间关系的图像序列。系统抽取帧、编码视觉内容，并保留帧顺序或时间信息，支持事件描述、动作比较与片段定位。视频指令数据中的问题与回答需要对应到正确的时间段，长视频任务还涉及帧采样和视觉信息压缩。

\textbf{语音输入。}语音编码器把声音转为可供模型使用的表示，后续可以连接语言模型完成转写、问答与对话。语音输出还需要声音生成模块或音频 token 解码器。实时交互进一步涉及流式输入、说话轮次和打断处理，训练数据相应包含声音与文本或声音与声音的配对。

图像、视频与语音的共同思路，是把各模态编码为可组合的表示，再通过任务数据学习如何读取和生成。本章介绍的视觉输入方案提供了理解多模态后训练的基本框架；下一章将进一步列出多模态生成、世界模型与具身智能等研究方向。

\textbf{本章相关资料}
\begin{enumerate}
    \item Radford 等，2021，\href{https://arxiv.org/abs/2103.00020}{Learning Transferable Visual Models From Natural Language Supervision}。阅读重点：CLIP 的图文双编码器、双向对比学习与零样本分类。
    \item Liu 等，2023，\href{https://arxiv.org/abs/2304.08485}{Visual Instruction Tuning}。阅读重点：视觉编码器、连接器与语言模型的组合，以及图文对齐和指令微调的分工。
    \item Bai 等，2025，\href{https://arxiv.org/abs/2502.13923}{Qwen2.5-VL Technical Report}。阅读重点：动态分辨率、视觉文本与文档理解，以及视频输入的处理方式。
\end{enumerate}
